\section{Cubic coincident Stieltjes products and a diagonal Tornheim jet}
\label{sec:cubic}

The incoming \emph{Coincident Stieltjes} article explicitly asks for the
unit-coordinate finite part of $\int_0^1\psi(x)^3\,dx$ and an all-index
three-factor closure theorem.  It suggests higher-depth coordinates and
does not assume a reduction to ordinary zeta constants.  The results
below supply an exact reduction to a specified diagonal Tornheim
derivative, together with a holomorphic germ for all mixed Stieltjes
indices.  They do not evaluate that remaining derivative in a
depth-one constant algebra.

The triple-Hurwitz/Tornheim correspondence is classical, as is the
Mellin--polylogarithm continuation used to compute Tornheim functions.
Relevant antecedents include Espinosa--Moll, Borwein--Dilcher,
Bailey--Borwein, and Onodera\cite{EMTornheim,BorweinDilcher,BaileyBorwein,Onodera2021}.
The low diagonal jets below are recovered identities from this established
setting.

\subsection{Normalization and the classical three-factor kernel}

Use the unit coordinate $x$ and, for $e\in\mathbb Z$ and
$k\in\mathbb Z_{\ge0}$, the convention
\[
 \operatorname{FP}\int_0^1 x^e\log^k x\,dx=
 \begin{cases}(-1)^k k!/(e+1)^{k+1},&e\ne-1,\\0,&e=-1.
 \end{cases}
\]
This agrees with the constant term after removing the divergent
endpoint powers and logarithms from $\int_\epsilon^1$.
Let
\[
 g(u,x)=\zeta(1+u,x)-\frac1u
       =\sum_{m\ge0}\frac{(-1)^m}{m!}\gamma_m(x)u^m,
 \qquad
 b(u)=\zeta(1+u)-\frac1u,
 \qquad
 a(u)=-(1+u)\zeta(2+u).
\]
The removable singularities in $g$ and $b$ are always removed.
In particular $g(0,x)=-\psi(x)$, $b(0)=\gamma$, and
$a(0)=-\zeta(2)$.

Initially in its absolutely convergent region, put
\[
 T(A,B,C)=\sum_{m,n\ge1}\frac1{m^A n^B(m+n)^C}.
\]
Its meromorphic continuation is understood in every later formula.
The initial absolute-convergence domain is given by
$\Re(A+C)>1$, $\Re(B+C)>1$, and $\Re(A+B+C)>2$.
Write $K_2(s,t)=\int_0^1\zeta(s,x)\zeta(t,x)\,dx$ and
$K_3(s_1,s_2,s_3)=\int_0^1\prod_{i=1}^3\zeta(s_i,x)\,dx$ in an
initial convergence domain, then continue meromorphically.
If $\lambda_i=1-s_i$, the classical triple kernel is
\begin{align}
 K_3(s_1,s_2,s_3)
 &=\frac{2\Gamma(\lambda_1)\Gamma(\lambda_2)\Gamma(\lambda_3)}
          {(2\pi)^{\lambda_1+\lambda_2+\lambda_3}}
 \sum_{k=1}^3
 \cos\frac{\pi(\lambda_i+\lambda_j-\lambda_k)}2\,
 T(\lambda_i,\lambda_j,\lambda_k),             \tag{C1}\label{cub:C1}
\end{align}
where $i<j$ are the indices other than $k$.
For $\Re\lambda_i>1$, the Fourier series
\[
 \zeta(1-\lambda,x)=\Gamma(\lambda)(2\pi)^{-\lambda}
 \sum_{n\ne0}|n|^{-\lambda}
       e^{-\pi i\lambda\operatorname{sgn}(n)/2}e^{2\pi inx}
\]
converges absolutely.  The zero-frequency condition in the product
requires two positive frequencies and one negative, or its opposite.
The two sign regions paired for a fixed negative position give the
cosine in (C1). This proves (C1) in a nonempty open set; meromorphic
continuation gives the stated identity. The same argument gives the pair kernel
\[
K_2(s,t)=\frac{2\Gamma(1-s)\Gamma(1-t)}{(2\pi)^{2-s-t}}
\cos\frac{\pi(s-t)}2\,\zeta(2-s-t).
\]

\subsection{The complete cubic subtraction}

For compactness write $K_{ij}=K_2(1+u_i,1+u_j)$,
$K_{123}=K_3(1+u_1,1+u_2,1+u_3)$, and $U=u_1+u_2+u_3$.
Define the meromorphic germ
\[
 I_3(\boldsymbol u)=K_{123}-\sum_{k=1}^3\frac{K_{ij}}{u_k}
                    -\frac1{u_1u_2u_3}.        \tag{C2}\label{cub:C2}
\]
The mean of a single Hurwitz factor is zero in the initial domain, so
this is precisely the meromorphic continuation of
$\int_0^1\prod_i g(u_i,x)\,dx$.

\begin{theorem}[Cubic holomorphic completion]
The germ
\[
 F_3(\boldsymbol u)=I_3(\boldsymbol u)
   +\sum_{k=1}^3\frac{a(u_k)}{u_i+u_j}
   +\sum_{k=1}^3\frac{b(u_i)b(u_j)}{u_k}       \tag{C3}\label{cub:C3}
\]
is holomorphic in a neighborhood of $\boldsymbol u=0$. For all
$m_1,m_2,m_3\ge0$,
\[
 \operatorname{FP}\int_0^1\prod_{i=1}^3\gamma_{m_i}(x)\,dx
 =(-1)^{m_1+m_2+m_3}\prod_{i=1}^3m_i!\,
   [u_1^{m_1}u_2^{m_2}u_3^{m_3}]F_3(\boldsymbol u). \tag{C4}\label{cub:C4}
\]
\end{theorem}

\begin{proof}
The shift relation gives
\[
 g(u,x)=x^{-1-u}+h(u,x),\qquad h(u,x)=g(u,1+x),
 \qquad h(u,x)=b(u)+a(u)x+O(x^2).
\]
All these Taylor statements are uniform on small compact spectral
neighborhoods.  Write $h_i=h(u_i,x)$.  The product of the three singular
terms contributes $-1/(2+U)$ after meromorphic integration. A term with
the two singular factors in positions $i,j$ and smooth factor $h_k$
contributes
\[
 -\frac{b(u_k)}{1+u_i+u_j}-\frac{a(u_k)}{u_i+u_j}
 +\int_0^1x^{-2-u_i-u_j}
             \bigl(h_k-b(u_k)-a(u_k)x\bigr)\,dx.
\]
A term with just the singular factor in position $k$ contributes
\[
 -\frac{b(u_i)b(u_j)}{u_k}
 +\int_0^1x^{-1-u_k}\bigl(h_i h_j-b(u_i)b(u_j)\bigr)\,dx.
\]
The remaining term is $\int_0^1h_1h_2h_3\,dx$.
After (C3), all remaining integrals are normally convergent near zero:
their endpoint majorants have the form
$Cx^{-\eta}(1+|\log x|^M)$, with $\eta<1$, after any prescribed finite
number of spectral derivatives.  This proves holomorphy.

The deleted expressions are exactly the integrations of resonant
monomials with exponent $-1$ at $\boldsymbol u=0$. Each coefficient of
such a monomial is a multiple of $x^{-1}\log^j x$, whose unit-coordinate
finite part is zero. The other integrated monomials have precisely the
finite parts prescribed above. Coefficient extraction under the
normally convergent remainder proves (C4).
\end{proof}

The mixed spectral derivatives in (C4) apply to the completed germ.
Uncompleted individual Tornheim summands can have intersecting polar
hyperplanes at the origin. In particular one must not read (C4) as a
formula involving unspecified ordinary partial derivatives of
$T$ at $(0,0,0)$.

\subsection{Every argument derivative}

The same proof gives a finite all-order extension. For
$p_i\in\mathbb Z_{\ge0}$, let
\begin{gather*}
 c_p(u)=(-1)^p(1+u)_p,\qquad
 h_i(u_i,x)=\partial_x^{p_i}g(u_i,1+x),\\
 G_i=\partial_x^{p_i}g(u_i,x)
     =c_{p_i}(u_i)x^{-p_i-1-u_i}+h_i(u_i,x).
\end{gather*}
Let $I_{\boldsymbol p}$ be the meromorphic product integral of the
$G_i$. With $s_i=1+p_i+u_i$, it is explicitly
\begin{align*}
 I_{\boldsymbol p}={}&\left(\prod_{i=1}^3 c_{p_i}(u_i)\right)
                         K_3(s_1,s_2,s_3)\\
 &-\sum_{k=1}^3\frac{\delta_{p_k0}c_{p_i}(u_i)c_{p_j}(u_j)}{u_k}
                         K_2(s_i,s_j)
 -\frac{\prod_{i=1}^3\delta_{p_i0}}{u_1u_2u_3}.
\end{align*}
The single-Hurwitz integrals vanish in the initial convergence domain
and hence identically after continuation.
Then the holomorphic completion is
\begin{align}
 F_{\boldsymbol p}=I_{\boldsymbol p}
 &+\sum_{k=1}^3\frac{c_{p_k}(u_k)}{u_k}
       [x^{p_k}]\bigl(h_i(u_i,x)h_j(u_j,x)\bigr)\notag\\
 &+\sum_{k=1}^3\frac{c_{p_i}(u_i)c_{p_j}(u_j)}{u_i+u_j}
       [x^{p_i+p_j+1}]h_k(u_k,x).             \tag{C5}\label{cub:C5}
\end{align}
Its coefficients generate
$\operatorname{FP}\int_0^1\prod_i\gamma_{m_i}^{(p_i)}(x)\,dx$ with
the same factorial and sign in (C4). Indeed a subset $S$ of singular
factors has exponent $-|S|-\sum_{i\in S}p_i-\sum_{i\in S}u_i$;
the unique resonant smooth Taylor degree is
$|S|+\sum_{i\in S}p_i-1$.  For a singleton or pair this gives precisely
(C5). For all three factors the needed smooth degree is positive but
there is no smooth factor, so no further resonant term occurs.
To justify the normal convergence rather than merely identify the
formal poles, subtract every smooth Taylor term up to that threshold.
The integrated lower Taylor degrees have nonzero integer denominators
at the origin, so they are holomorphic and retain their prescribed
monomial finite parts. The only denominators vanishing at the origin
are exactly those displayed in (C5). The remaining Taylor remainders
are integrable uniformly on sufficiently small compact spectral
neighborhoods, also after any fixed finite number of derivatives.
Every smooth coefficient is a finite combination of Hurwitz jets,
because
\[
 [x^\ell]h_i(u_i,x)
 =\frac{c_{p_i+\ell}(u_i)}{\ell!}
      \zeta(1+p_i+\ell+u_i)
   -\frac{\delta_{p_i0}\delta_{\ell0}}{u_i}.
\]
The Pochhammer coefficient expansions are the usual finite harmonic
polynomials. Thus (C5) is an explicit finite recipe for all cubic
argument-derivative moments.

\subsection{An independently defined diagonal Tornheim function}

Let $\tau(s)=T(s,s,s)$ initially for $\Re s>2/3$, and then continue
this one-variable function.  The diagonal restriction is essential:
$T$ itself has no direction-independent value at the multivariate
origin. For example, $T(s,s,0)=\zeta(s)^2$ tends to $1/4$, whereas
the diagonal limit is $1/3$, as recovered in (C8) below.

Here is a convergent formula defining the germ of $\tau$ near zero
without using any Hurwitz product integral. Set
$L_s(t)=\operatorname{Li}_s(e^{-t})$, $G(s)=\Gamma(1-s)$, and
\begin{align*}
 R_s(t)={}&L_s(t)^2-G(s)^2t^{2s-2}
   -2G(s)\zeta(s)t^{s-1}
   +2G(s)\zeta(s-1)t^s-\zeta(s)^2,\\
 E(s)={}&\int_0^1t^{s-1}R_s(t)\,dt
          +\int_1^\infty t^{s-1}L_s(t)^2\,dt.
\end{align*}
For sufficiently small $\delta>0$, $E$ is holomorphic when $|s|<\delta$.
Indeed the Jonqui\`ere expansion \cite{DLMFPolylog}
\[
 L_s(t)=\Gamma(1-s)t^{s-1}
        +\sum_{k=0}^\infty\frac{(-1)^k}{k!}\zeta(s-k)t^k,
 \qquad |t|<2\pi,
\]
shows that the subtracted integrand is bounded by
$C(t^{-2\delta}+t^{-\delta})(1+|\log t|^M)$ after any fixed number
of spectral derivatives. The second integral converges exponentially.
Direct Mellin integration first for $2/3<\Re s<1$ and then continuation
gives the holomorphic formula
\[
 \tau(s)=\frac1{\Gamma(1+s)}\left\{
 sE(s)+\frac{s\Gamma(1-s)^2}{3s-2}
 +\frac{2s\Gamma(1-s)\zeta(s)}{2s-1}
 -\Gamma(1-s)\zeta(s-1)+\zeta(s)^2\right\}. \tag{C6}\label{cub:C6}
\]
This is the $\theta=1$ local-subtraction form of the classical
Crandall Mellin continuation. In particular $\tau'''(0)$ is an
unambiguous, independently computable spectral constant.

\subsection{The cubic digamma reduction}

Put $\ell=\log(2\pi)$, and define
\[
 Z(t)=t\zeta(1+t),\quad
 A_d(t)=\frac{\Gamma(1-t)^2}{\Gamma(1-2t)\cos(\pi t)},\quad
 B(t)=\Gamma(1-t)^3(2\pi)^{3t}\cos(\pi t/2).
\]
The classical pair formula yields
$K_2(1+t,1+t)=-A_d(t)Z(2t)/t^2$, while (C1) gives
$K_3(1+t,1+t,1+t)=-6B(t)\tau(-t)/t^3$. Hence
\[
 F_3(t,t,t)=
 \frac{-6B(t)\tau(-t)+3A_d(t)Z(2t)-1+3(Z(t)-1)^2}{t^3}
 +\frac{3a(t)}{2t}.                           \tag{C7}\label{cub:C7}
\]
Since the left side is holomorphic, comparison of the three negative
powers gives
\[
 \tau(0)=\frac13,\qquad
 \tau'(0)=\ell,\qquad
 \tau''(0)=3\ell^2-2\gamma^2-4\gamma_1+\zeta(2).\tag{C8}\label{cub:C8}
\]
These are recovered low-jet identities; the first two are well known,
and Onodera's 2021 paper evaluates second derivatives at the origin.

For completeness the needed expansions are
\begin{align*}
 Z(t)&=1+\gamma t-\gamma_1t^2+\tfrac12\gamma_2t^3+O(t^4),\\
 A_d(t)&=1+2\zeta(2)t^2-2\zeta(3)t^3+O(t^4),\\
 \log B(t)&=3(\gamma+\ell)t+\tfrac34\zeta(2)t^2
                 +\zeta(3)t^3+O(t^4),\\
 a'(0)&=-\zeta(2)-\zeta'(2).
\end{align*}
Taking the constant coefficient of (C7) and using
$F_3(0,0,0)=-\operatorname{FP}\int_0^1\psi(x)^3\,dx$ proves
\begin{align}
 \operatorname{FP}\int_0^1\psi(x)^3\,dx={}&-\tau'''(0)
 -9\gamma^3-18\gamma^2\ell-30\gamma\gamma_1
 -36\ell\gamma_1-12\gamma_2+9\ell^3\notag\\
 &+\frac32\gamma\zeta(2)+9\ell\zeta(2)
 +\frac32\zeta(2)+8\zeta(3)+\frac32\zeta'(2). \tag{C9}\label{cub:C9}
\end{align}
An equivalent unsimplified coefficient certificate is
\begin{align*}
 F_3(0,0,0)={}&-6[t^3]B(t)\tau(-t)
 +12\gamma_2+12\gamma\zeta(2)-6\zeta(3)-6\gamma\gamma_1\\
 &-\tfrac32\bigl(\zeta(2)+\zeta'(2)\bigr).
\end{align*}

\subsection{Directional Laurent constants do not fix the coordinate finite part}

Let $\alpha_1,\alpha_2,\alpha_3$ be fixed complex numbers with
$\alpha_i\ne0$ and $\alpha_i+\alpha_j\ne0$. If $\operatorname{CT}$
denotes the one-variable Laurent constant and
$Q=\operatorname{FP}\int_0^1\gamma_0(x)^3\,dx$, then (C3) gives
\begin{align}
 Q={}&\operatorname{CT}_{t=0} I_3(\alpha_1t,\alpha_2t,\alpha_3t)
 -\bigl(\zeta(2)+\zeta'(2)\bigr)
       \sum_{k=1}^3\frac{\alpha_k}{\alpha_i+\alpha_j}\notag\\
 &-\gamma\gamma_1\sum_{k=1}^3\frac{\alpha_i+\alpha_j}{\alpha_k}.
                                                        \tag{C10}\label{cub:C10}
\end{align}
For the diagonal this correction is
$-\frac32(\zeta(2)+\zeta'(2))-6\gamma\gamma_1$.
This is a normalization warning, not an identified false claim in an
inspected source. Removing only the negative Laurent powers after
diagonal specialization would omit the finite coefficients of the
resonant expressions. The full germ (C3) removes exactly those entire
expressions and yields the fixed-coordinate finite part.

\subsection{Independent numerical diagnostic}

The script \texttt{verify\_cubic.py} computes the first four diagonal
Tornheim derivatives from (C6), using the convergent Jonqui\`ere series
on $(0,1)$ and an exponentially convergent integral on $(1,\infty)$.
The independent digamma calculation integrates
\[
 \psi(x)^3+x^{-3}+3\gamma x^{-2}
                -3\bigl(\zeta(2)-\gamma^2\bigr)x^{-1}
\]
ordinarily and then adds $1/2+3\gamma$. Near zero it evaluates the
regular remainder by the convergent Taylor series for $\psi(1+x)$,
avoiding cancellation of huge singular terms.
At 35 working decimal digits and 50 Jonqui\`ere terms the values are
\begin{align*}
 \tau'''(0)&\approx86.657866601100073652903276858762183,\\
 \operatorname{FP}\int_0^1\psi(x)^3\,dx
 &\approx1.9662627343915343406643753384116314.
\end{align*}
The two independent expressions in (C9) differ by less than
$4.1\times10^{-34}$. This is a floating-point diagnostic, not an
interval enclosure or substitute for the proof above.

